
\def\PUDcodigo{EME114}

\def\PUDcodacad{04506.27}

\def\PUDdisciplina{Termodinâmica}

\def\PUDsemestre{S6}

\def\PUDhoras{80}

%NOVOS ---------------------------------------------------
\def\PUDhorasTeoria{76}
\def\PUDhorasPratica{4}

\def\PUDinterECA{
\hyperref[PUD:EME115]{Mecanismos (04506.25)}

\hyperref[PUD:EME119]{Sistemas Mecânicos (04506.35)}
}

\def\PUDatuaisECA{---}

\def\PUDrecursos{Projetor multimídia; Quadro branco e pincel; Aulas em laboratório.}

%----------------------------------------------------------

\def\PUDcreditos{4}

\def\PUDprerequisito{ \hyperref[PUD:ACE010]{ACE010}  \hyperref[PUD:ACE011]{ACE011} }

\def\PUDpreacad{ \hyperref[PUD:ACE010]{Cálculo 3 (04506.12)} 

\hyperref[PUD:ACE011]{Física 2 (04506.13)} 
}

\def\PUDobjetivos{Familiarizar-se com a termodinâmica clássica, como forma de ter base para subsequentes estudos em áreas como mecânica dos fluidos e máquinas térmicas; Ser capaz de fazer uso efetivo da termodinâmica na prática da engenharia; Compreender os fenômenos relativos à mudança de estados; Compreender os fenômenos relativos à conservação da energia; Compreender os fenômenos relativos às irreversibilidades nos processos termodinâmicos.}

\def\PUDementa{Propriedades de uma substância pura;  Trabalho e calor;  Primeira Lei da Termodinâmica;  Segunda Lei da Termodinâmica;  Entropia;  Irreversibilidade e Disponibilidade.}

\def\PUDconteudo{ & Alguns comentários preliminares: Central termoelétrica, Célula de combustível, Ciclo de refrigeração por compressão a vapor, Refrigerador termoelétrico, Equipamento de decomposição do ar, Turbina a gás, Motor químico de foguete. 
\\ 
 & Alguns conceitos e definições: O sistema termodinâmico e o volume de controle, Pontos de vista macroscópico e microscópico, Estado e propriedades de uma substância, Processos e ciclos, Unidades de massa, comprimento, tempo e força, Energia, Volume específico e massa específica, Pressão, Igualdade de temperatura, A lei zero da termodinâmica, Escalas de temperatura, Aplicações na Engenharia. 
\\ 
 & Propriedades de uma substância pura: A substância pura, Equilíbrio entre fases vapor-líquida-sólida para uma substância pura, Propriedades independentes de uma substância pura, Tabelas de propriedades termodinâmicas, Superfícies termodinâmicas, O comportamento P-V-T dos gases na região de massas específicas pequenas ou moderadas, O fator de compressibilidade, Equações de estado, Aplicações na Engenharia. 
\\ 
 & Trabalho e calor: Definição de trabalho, Unidades de trabalho, Trabalho realizado na fronteira móvel de um sistema simples compressível, Considerações finais sobre trabalho, Definição de calor, Modos de transferência de calor, Comparação entre calor e trabalho, Aplicações na engenharia. 
\\ 
 & Primeira Lei da Termodinâmica: A primeira lei da termodinâmica para um sistema que percorre um ciclo, A primeira lei da termodinâmica para uma mudança de estado num sistema, Energia interna – uma propriedade termodinâmica, A propriedade termodinâmica entalpia, Calores específicos a volume e a pressão constantes, Energia interna, Entalpia, e Calor específico de gases ideais, Equação da primeira lei em termos de taxas, Conservação da massa, Aplicações na engenharia. 
\\ 
 & Segunda Lei da Termodinâmica: Motores térmicos e refrigeradores, Segunda lei da termodinâmica, O processo reversível, Fatores que tornam um processo irreversível, O ciclo de Carnot, A escala termodinâmica de temperatura, A escala de temperatura de gás ideal, Máquinas reais e ideais, Aplicações na engenharia. 
\\ 
 & Entropia: Desigualdade de Clausius, Entropia - uma propriedade do sistema, A entropia para uma substância pura, Variação de entropia em processos reversíveis, duas relações termodinâmicas importantes, Variação de entropia num sólido ou líquido, Variação de entropia num gás ideal, Processo politrópico reversível para um gás ideal, Variação de entropia do sistema durante um processo irreversível, geração de entropia, princípio de aumento de entropia, Equações da taxa de variação de entropia, Comentários gerais sobre entropia e caos. 
\\ 
 & Irreversibilidade e Disponibilidade: Energia disponível, Trabalho reversível e Irreversibilidade, Disponibilidade e eficiência baseada na segunda lei da termodinâmica, Equação de balanço de exergia, Aplicações na engenharia.
 }

\def\PUDmetodo{Aulas expositórias que podem ser teóricas e/ou práticas, onde as práticas no laboratório serão marcadas no decorrer da disciplina. Um pequeno projeto experimental que transforme calor em trabalho é passado logo no início da disciplina.}

\def\PUDavalia{A avaliação será feita com aplicação de provas teóricas e/ou práticas, além da possibilidade de inclusão de trabalhos e seminários no decorrer da disciplina. Conta como avaliação o memorial de cálculo de um pequeno projeto termodinâmico que transforma calor em trabalho.
}

\def\PUDbibbasica{ & \href{www.biblioteca.ifce.edu.br/index.asp?codigo_sophia=65410}{Borgnakke}, C.; Sonntag, R. E.  Fundamentos da Termodinâmica.  São Paulo, SP: Edgard Blücher, 2013.  728p.  ISBN: 9788521207924. 
\\ 
 &  \href{www.biblioteca.ifce.edu.br/index.asp?codigo_sophia=24058}{Moran}, M. J.;  Shapiro, H. N.;   Muson, B. R.;   DeWitt, D. P.  Introdução à engenharia de sistemas térmicos: termodinâmica, mecânica dos fluidos e transferência de calor.  Rio de Janeiro, RJ: LTC, 2005.  604p.  ISBN: 9788521614463.  
\\ 
 &  \href{www.biblioteca.ifce.edu.br/index.asp?codigo_sophia=63341}{Moran}, M. J.;  Shapiro, H. N.;  Princípios de Termodinâmica para Engenharia.  7ª ed.  Rio de Janeiro, RJ: LTC, 2016.  819p.  ISBN: 9788521622123. }

\def\PUDbibcomplementar{ & \href{www.biblioteca.ifce.edu.br/index.asp?codigo_sophia=41355}{Halliday}, D.;  Resnick, R.;  Walker, J.  Fundamentos de física: gravitação, ondas e termodinâmica.  Vol. 2.  9ª ed.  Rio de Janeiro, RJ: LTC, 2013.  296p.  ISBN: 9788521619048. 
\\ 
 &  \href{www.biblioteca.ifce.edu.br/index.asp?codigo_sophia=63133}{Fox}, R.  W.;  MacDolnad, A.  T.;   Pritchard, P.  J.  Introdução à Mecânica dos Fluidos.  8ª ed. Rio de Janeiro, RJ: LTC, 2014.  871p.  ISBN: 9788521623021. 
\\ 
 &  \href{www.biblioteca.ifce.edu.br/index.asp?codigo_sophia=61984}{Incropera}, F. P.;   DeWitt, D. P.;   Bergman, T. L.;   Lavine, A. S.  Fundamentos de Transferência de Calor e de Massa.  7ª ed. Rio de Janeiro, RJ: LTC, 2014.  672p.  ISBN: 9788521625049.
\\
 &  \href{www.biblioteca.ifce.edu.br/index.asp?codigo_sophia=40619}{Livi}, Celso Pohlmann.  Fundamentos de fenômenos de transporte: um texto para cursos básicos.  2º ed.  Rio de Janeiro, RJ: LTC, 2012.  237p.  ISBN: 9788521620570.
\\
 &  \href{www.biblioteca.ifce.edu.br/index.asp?codigo_sophia=41209}{Azevedo}, Edmundo J. S. Gomes de.  Termodinâmica aplicada.  3º ed.  Lisboa (Portugal): Escolar, 2011.  891p.  ISBN: 9789725923153. }

\def\PUDautor{Francisco Frederico dos Santos Matos}

\def\PUDdata{2013-04-25}

\def\PUDrevisao{2}

\def\PUDrevisor{Francisco Frederico dos Santos Matos}

\def\PUDdatarevisao{2019-05-15}

\PUDcorpo
